\section{Random routes and encounters away from a common endpoint}\label{rt:section}\label{sec:routes}

A transport between two deletion measures follows one route.  Its second
moment compares two independent route choices conditional on the same
exposed datum and first tag, with the two other tags ranging over the
distant cell.  The two routes therefore have one endpoint in common, but their
other endpoints need not coincide.  We construct their geometric law
here, before estimating the transport likelihoods in
Section~\ref{long:section}.

Two properties will be needed.  The total number of encounters has a
bounded exponential moment.  If an initial portion of each route is
discarded, the probability of any remaining encounter tends to zero
when the other endpoints are averaged in a growing cube.  We call the
discarded portion \emph{silent}: the corresponding trajectories will be
kept unchanged in the observed copy of the transport.  These estimates
include the average over the random set of unusable sites; they are not
estimates for each fixed configuration of that set.

We adapt the folded-path construction of \cite[Section~7]{H}.  Its
probability estimate is an instance of the unpredictable-path method of
Benjamini, Pemantle, and Peres~\cite{BPP}.  We give the finite process
and the estimates needed for unequal terminal endpoints explicitly.

Write \(\mathcal L_K=(\mathbb Z/K\mathbb Z)^3\), and let
\(d_\infty\) be its periodic maximum distance.  Sites differing by one
in one coordinate are nearest neighbors.  For an integer \(a\ge1\),
range-\(a\) adjacency means maximum distance at most \(a\).
A random datum \(\mathcal D\) determines a bad set
\(\mathcal B\subset\mathcal L_K\); its complement consists of the
good sites.  The sole probabilistic assumption
on this set is
\begin{equation}\label{rt:bad-law}
 \mathbb P(A\subset\mathcal B)\le p^{|A|}
 \qquad(A\subset\mathcal L_K).
\end{equation}
No independence or translation invariance is assumed.

Here and below a cube of side at most \(c_QR\) on the torus is the
projection of an ordinary lattice cube of that size; \(c_Q\) is fixed.
We take \(K\ge C_0R\), with \(C_0\) sufficiently large in terms of
the fixed constants.  Fix a common endpoint \(z\) and such a cube \(Q\)
so that every \(y\in Q\) has forward first-coordinate displacement
\(n_y\in[.29K,.46K]\) from \(z\).  Rounding the endpoints of this
interval has no effect on the estimates.  Fix an integer contact radius
\(r_c\ge1\), and a nonnegative integer \(q\), which allows a fixed
number of neighbors of each changed portion to be included in the
contact count.  All constants in this section may depend on
\(r_c,q,c_Q\), but not on \(K,R\).

\begin{proposition}[Routes with rare encounters away from the common end]
\label{rt:routes}
There are \(p_0,\kappa,c,C>0\) and \(R_0\) with the following
property.  Suppose \(R\ge R_0\), \(K\ge C_0R\), and
\eqref{rt:bad-law} holds with \(p\le p_0\).  Put
\(h=\lfloor R^{1/2}\rfloor\).  There is an event \(\mathcal E\),
determined by \(\mathcal D,z,Q\), on which the construction below gives
simple good nearest-neighbor routes \(\pi_y\) from \(z\) to every
\(y\in Q\).  Conditional on the datum, routes are sampled using fresh
independent random variables.  The event includes the requirement that
the \(10R\)-neighborhoods of \(z\) and \(Q\) contain no bad sites.

For a route \(\pi=(v_0,\ldots,v_m)\) starting at \(z\), set
\[
 a(\pi)=\max\{i:d_\infty(v_i,z)\le h\},\qquad
 \pi^{\rm ch}=\{v_i:i>a(\pi)\}.
\]
Its silent portion \((v_0,\ldots,v_{a(\pi)})\) has at most \(Ch\)
vertices.  The analogous silent portion read backwards from \(y\)
also has at most \(Ch\) vertices, and the two portions are disjoint.
Let \(\pi^{\rm ch,q}\) consist of vertices whose path index is within
\(q\) of an index in \(\pi^{\rm ch}\).  For independently sampled
routes \(\pi_y,\pi_{y'}\), define
\begin{align*}
 J&=\sum_{v\in\pi_y}\sum_{v'\in\pi_{y'}}
       \mathbf1_{\{d_\infty(v,v')\le r_c\}},\\
 J_*&=\sum_{v\in\pi_y^{\rm ch,q}}
       \sum_{v'\in\pi_{y'}^{\rm ch,q}}
       \mathbf1_{\{d_\infty(v,v')\le r_c\}}.
\end{align*}
For every fixed pair \(y,y'\in Q\),
\begin{equation}\label{rt:exponential}
 \mathbb E_{\mathcal D}\!\left[
  \mathbf1_{\mathcal E}\,
  \mathbb E_{\rm routes}(e^{\kappa J}\mid\mathcal D)\right]\le C.
\end{equation}
For every deterministic probability mass function \(\omega\) on
\(Q\times Q\) satisfying \(\omega(y,y')\le C_\omega R^{-6}\),
\begin{equation}\label{rt:rare}
 \sum_{y,y'\in Q}\omega(y,y')\,
 \mathbb E_{\mathcal D}\!\left[
  \mathbf1_{\mathcal E}\,
  \mathbb P_{\rm routes}(J_*>0\mid\mathcal D)\right]
 \le CC_\omega R^{-c_g},\qquad
 c_g=\frac12\left(\frac{\log2}{\log3}-\frac12\right)>0.
\end{equation}
The same conclusions hold with a common terminal endpoint, reading the
routes in reverse order.

If \(\mathcal M\) denotes just the stated bad-free margin event, then
\begin{equation}\label{rt:rejection}
 \mathbb P(\mathcal M\setminus\mathcal E)
 \le K^3(Ap)^{\lfloor K/(10(r_c+2))\rfloor}
       +CR^3e^{-cR}.
\end{equation}
Here \(A\) is fixed and \(p_0\) is chosen so that \(Ap_0<1\).
\end{proposition}

\subsection{A finite process with small conditional atoms}

We first produce the transverse randomness.  Its conditional atom bound
will also control encounters with an arbitrary fixed opposing route.

\begin{lemma}\label{rt:prediction}
Put \(\sigma=\log2/\log3\).  For every integer \(T\ge1\) there
is an integer-valued process \((S_j)_{j=0}^T\), with \(S_0=0\) and
\(|S_{j+1}-S_j|=1\), such that
\begin{equation}\label{rt:prediction-bound}
 \sup_{x\in\mathbb Z}
 \mathbb P(S_{u+k}=x\mid S_0,\ldots,S_u)
 \le 6^\sigma k^{-\sigma}
 \quad(k\ge1,\ u+k\le T).
\end{equation}
\end{lemma}
\begin{proof}
Choose an ordered ternary tree of depth \(m\) with \(3^m\ge T\).
Give the root spin \(+1\).  At every nonleaf vertex the first two
children inherit its spin and the third receives an independent fair
spin.  If the leaf spins in order are \(X_1,\ldots,X_{3^m}\), let
\(S_j=\sum_{i=1}^jX_i\).

The leaf sum of a depth-\(d\) subtree with root spin \(s\) equals
\[
 2^ds+\sum_{a=1}^d2^{d-a}
                 \sum_{i=1}^{3^{a-1}}\varepsilon_{a,i},
\]
where the displayed fresh spins are independent.  Select one fresh spin
at each depth.  Their weighted sum is uniform on \(2^d\) distinct
integers: writing \(\varepsilon=2\eta-1\) reduces this assertion to
binary expansion.  Independently adding the remaining terms shows that
every atom of the subtree sum is at most \(2^{-d}\).

For \(k\ge2\), choose \(d\ge0\) with
\(2\cdot3^d\le k<2\cdot3^{d+1}\).  The interval of leaves
\(\{u+1,\ldots,u+k\}\) contains a complete aligned depth-\(d\)
subtree.  Condition on every fresh spin except those strictly below its
root.  This fixes the history through time \(u\), its root spin, and
all leaves outside the subtree.  The conditional atom bound just proved
is \(2^{-d}\le6^\sigma k^{-\sigma}\).  Conditioning down to the
history gives \eqref{rt:prediction-bound}.  The bound for \(k=1\) is
immediate.
\end{proof}

Take independent copies for the two transverse coordinates and write
\(X_s=(S_s^{(2)},S_s^{(3)})\).  Let \(\mathcal F_s\) be their joint
history.  For every deterministic transverse target \(a\), integer
radius \(l\ge0\), and \(k\ge1\), independence gives
\begin{equation}\label{rt:target-atoms}
 \mathbb P(\|X_{u+k}-a\|_\infty\le l\mid\mathcal F_u)
 \le C(1+l)^2k^{-2\sigma}.
\end{equation}
The same estimate holds for periodic distance modulo \(K\), provided
the horizon is at most \(K/2\): the possible values of each coordinate
lie in an interval of length at most \(K\), so only a bounded number
of representatives of a periodic target need be counted.  The bound
is also valid for a union of any fixed number of target squares, after
changing its constant.

\begin{lemma}[Ordered tests against moving targets]\label{rt:ordered}
At each time \(s\in\{0,\ldots,T\}\), let \(A_s(l)\) be a union of
at most two deterministic squares of maximum radius \(l\), with
arbitrary centers, in the transverse torus.  There is a constant
\(C_1\), independent of their centers and of \(T\le K/2\), such that
\begin{equation}\label{rt:ordered-bound}
 \sum_{0\le s_1<\cdots<s_j\le T}
 \mathbb P\{X_{s_i}\in A_{s_i}(l_i),\ 1\le i\le j\}
 \le \prod_{i=1}^jC_1(1+l_i)^2.
\end{equation}
In particular, for \(a_l\ge0\), if
\(C_1\sum_{l\ge0}a_l(1+l)^2<1\), then
\begin{equation}\label{rt:ordered-product}
 \mathbb E\prod_{s=0}^T
 \left(1+\sum_{l\ge0}a_l\mathbf1_{\{X_s\in A_s(l)\}}\right)
 \le \left(1-C_1\sum_{l\ge0}a_l(1+l)^2\right)^{-1}.
\end{equation}
\end{lemma}
\begin{proof}
Set \(k(0)=1\) and \(k(t)=t^{-2\sigma}\) for \(t\ge1\).
Successive conditioning in \eqref{rt:target-atoms} bounds a summand
on the left of \eqref{rt:ordered-bound} by
\[
 C^j\prod_{i=1}^j(1+l_i)^2\,
 k(s_1)\prod_{i=2}^j k(s_i-s_{i-1}).
\]
Since \(2\sigma>1\), the sum of \(k\) is finite.  Summing over the
initial time and positive gaps proves the first assertion.  Expanding
the product into ordered subsets of times and using nonnegative
summation proves the second.
\end{proof}

\subsection{Folding the walk and avoiding bad components}

For each \(y\in Q\), choose lifts of \(z,y\) with first-coordinate
difference \(n=n_y\) and transverse differences
\(\Delta_2,\Delta_3\in[-K/2,K/2]\).  Independently of the datum,
sample two copies of Lemma~\ref{rt:prediction}, through time
\(\lfloor n/2\rfloor\), and set
\begin{equation}\label{rt:nominal}
 u_i=z+\left(i,
       \left\lfloor\frac{i\Delta_2}{n}\right\rfloor
                         +S^{(2)}_{s_n(i)},
       \left\lfloor\frac{i\Delta_3}{n}\right\rfloor
                         +S^{(3)}_{s_n(i)}\right),\qquad
 s_n(i)=\min(i,n-i),\quad 0\le i\le n.
\end{equation}
We use the same notation for a lifted vertex and its projection when
the distance in question makes the interpretation explicit.  The
endpoints are \(u_0=z,u_n=y\).  Each coordinate changes by at most
three per step, since \(|\Delta_a|/n\le50/29<2\).  Refine a step by
changing its coordinates in a fixed order.  It has at most seven edges,
and all its vertices lie within maximum distance three of its initial
nominal vertex.  Call the resulting walk \(\Gamma_y\); repeated
vertex occurrences are retained at this stage.

We recall precisely the finite-component facts used to modify this
walk.  If \(C\subset\mathbb Z^3\) is finite and range-one connected,
its fill is \(C\) together with all finite nearest-neighbor components
of \(\mathbb Z^3\setminus C\).  Its fill lies in the coordinate
bounding box of \(C\): a vertex outside the box has a coordinate ray
to infinity avoiding \(C\).  The exterior shell
\[
 \partial_*^{\rm ext}C
 =\{v\notin\operatorname{fill}(C):d_\infty(v,C)=1\}
\]
is nearest-neighbor connected and has at most \(26|C|\) vertices.
The connectivity assertion is the finite exterior-boundary theorem in
\cite[Theorem~4]{Timar}; this is the boundary convention used in
\cite[Section~7]{H}.  Indeed a vertex outside the fill belongs to the
infinite nearest-neighbor component of the complement.  The displayed
shell is therefore precisely the outer star-boundary visible by a
nearest-neighbor path from infinity, as in that theorem.  If \(C\) is a maximal
range-one bad component, every vertex in its exterior shell is good.

For completeness, small components on the torus have unambiguous
finite lifts.  A range-\(a\) component of size \(m\), with
\(am<K\) and \(K>2a\), can be lifted by a spanning tree, using the
unique short displacement in \([-a,a]^3\) on every edge.  All lifted
vertex differences have maximum norm at most \(a(m-1)\).  On a
remaining edge, the discrepancy between this difference and its short
displacement is a multiple of \(K\), with norm at most \(am<K\),
and hence is zero.  Thus the full inverse image is a disjoint union of
finite translates, each projecting bijectively onto the component.

There are at most \(A_a^m\) range-\(a\) connected sets of size
\(m\) containing a specified root.  Indeed choose a deterministic
spanning tree and record its depth-first traversal, which has
\(2(m-1)\) steps from a fixed finite set.  The visited vertices recover
the set.  Consequently,
\begin{equation}\label{rt:animal}
 \mathbb P\{\text{a range-\(a\) bad component has size at least }m\}
 \le K^3(A_ap)^m.
\end{equation}
A larger component contains a connected subset of exactly \(m\)
sites, so the same estimate covers all larger sizes.

Set \(a=r_c+2\) and \(m_K=\lfloor K/(10a)\rfloor\).  Define
\(\mathcal E\) to be the bad-free margin event \(\mathcal M\)
intersected with the following two events: no range-\(a\) bad component
has \(m_K\) or more sites; and none of the sites in \(\{z\}\cup Q\)
belongs to the projection of a filled range-one bad component.  The
finite lifts just proved make this definition precise.

On \(\mathcal E\), replace visits of the lifted refinement to each
bad range-one component \(C\) as follows.  The vertices immediately
before its first visit and after its last visit lie in its exterior
shell.  The initial segment to the first such vertex avoids \(C\)
and connects to an endpoint outside its fill; the final segment has the
same property.  Join these two shell vertices by a simple
nearest-neighbor shell path and replace the intervening walk segment.
Each replacement removes a bad visit and introduces none.  Every bad
component subsequently treated was therefore met by the original
refinement.  When no bad visits remain, project the walk onto the torus
and erase loops in chronological order.  This gives \(\pi_y\).
All choices can be made by fixed orderings on finite sets, so the route
is a measurable function of the datum and its independent nominal
randomness.

To prove \eqref{rt:rejection}, the first rejected event is bounded by
\eqref{rt:animal}.  If a filled range-one component of size \(m\)
contains a given endpoint, its bounding box contains that endpoint.
It therefore has a root within distance \(m\), giving at most
\(Cm^3\) possible roots.  On \(\mathcal M\), all its sites are at
distance more than \(10R\) from that endpoint, so \(m\ge cR\).
Summing the connected-set bound over \(m\ge cR\) and over at most
\(CR^3\) candidate endpoints gives \(CR^3e^{-cR}\), after decreasing
\(p_0\).  This proves the claimed rejection estimate.

The silent portions require a separate deterministic observation.
No inserted shell vertex can lie within distance \(5h\) of either
endpoint when \(R\) is sufficiently large: every shell vertex has a
bad neighbor, whereas the \(10R\)-margin is bad-free.  A refinement
vertex within distance \(h\) of \(z\) belongs to a nominal step with
index at most \(h+3\).  This uses \(0\le i\le .46K<K/2\), so its
first-coordinate distance to \(z\) does not wrap around the torus.
The entire original refinement before that step has at most
\(7(h+4)\) edges and stays in the bad-free margin.  Detour replacement
and chronological loop erasure can delete vertices of this initial
segment but cannot insert vertices into it: a detour begins next to
a bad visit, and all such visits occur after this segment.  Thus every
surviving visit to the \(h\)-ball has at most \(Ch\) predecessors in
the final route.  At the other end, every occurrence of a nominal
refinement vertex within distance \(h\) of \(y\) lies in the last
\(h+4\) nominal steps.  That whole terminal segment is bad-free.
Detours introduce no vertex into it, and loop erasure preserves the
order of surviving occurrences, so a surviving such vertex has at most
\(Ch\) successors.  This proves the terminal assertion without any
reversibility assumption on chronological loop erasure.  The two
original end segments lie in disjoint
\(Ch\)-neighborhoods, and cannot overlap in path order because their
first-coordinate indices are respectively at most \(h+3\) and at
least \(n-h-3\).  Their surviving silent portions are disjoint as well.
Finally, every vertex of \(\pi^{\rm ch,q}\) has distance greater
than \(h-q\) from \(z\), since nearest-neighbor edges change this
distance by at most one.


\subsection{Exponential moments for unequal terminal endpoints}

We now estimate encounters without assuming that the two terminal
endpoints agree.  Write \(u_i\), \(0\le i\le n\), and
\(u'_j\), \(0\le j\le n'\), for two independent nominal walks
starting at \(z\), and let \(\Gamma,\Gamma'\) be their refinements.
Set
\begin{equation}\label{rt:clamp}
 t(i)=\min(i,n'),\qquad d_i=d_\infty(u_i,u'_{t(i)}).
\end{equation}
Clamping to the last available opposing index is needed when
\(n\ne n'\).  The lifted first coordinates give
\(|i-j|=d_K(i,j)\), since both indices lie in \([0,.46K]\).
If two refined vertices from steps \(i,j\) approach within radius
\(l\), then \(|i-j|\le l+6\).  The map \(t\) is
one-Lipschitz and \(|j-t(i)|\le|j-i|\).  Bounded speed therefore
gives
\begin{equation}\label{rt:alignment}
 d_\infty(u_i,u'_j)\le l
 \quad\Longrightarrow\quad d_i\le C(l+1).
\end{equation}
The same implication, with adjusted constants, holds when refined
vertices replace the nominal vertices.  Since a step has a bounded
number of vertex occurrences and only boundedly many opposing indices
can contribute at a fixed radius, nominal encounters satisfy
\begin{equation}\label{rt:nominal-charge}
 J_{r_c}(\Gamma,\Gamma')
 \le C\sum_{i=0}^n\mathbf1_{\{d_i\le L\}}
\end{equation}
for a fixed \(L\).  The contact count on refinements counts vertex
occurrences, which only enlarges this bound for their surviving
vertices.

Let \(F\) range over the actual range-\(a\) bad components, with
\(a=r_c+2\).  Every encounter involving a shell vertex can be charged
to a component \(F\) satisfying
\(d_\infty(F,\Gamma)\le a\) and
\(d_\infty(F,\Gamma')\le a\).  Indeed a shell vertex is adjacent
to its generating range-one component, which was met by its own
original refinement.  An opposing nominal vertex within distance
\(r_c\) puts that component within distance \(a\) of the other
refinement.  If both vertices lie in shells of components \(C,C'\),
then \(d_\infty(C,C')\le r_c+2=a\); thus \(C,C'\) belong to
one range-\(a\) component.  There are at most \(26|F|\) possible
shell sites associated with \(F\), and each meets at most
\((2r_c+1)^3\) vertices of the other simple route.  Consequently,
on \(\mathcal E\),
\begin{equation}\label{rt:detour-charge}
 J\le J_{r_c}(\Gamma,\Gamma')
       +C\sum_F |F|\,
       \mathbf1_{\{d_\infty(F,\Gamma)\le a,
                      d_\infty(F,\Gamma')\le a\}}.
\end{equation}
Simplicity of the final routes is used here; no length bound on the
original refinements is used in charging a shell site.

Hold both complete nominal walks fixed and average over the environment
first.  Expand the exponential of the component sum in
\eqref{rt:detour-charge} as a product over actual components.  Every
term selects disjoint connected sets.  By \eqref{rt:bad-law}, its
probability is at most \(p\) to the sum of their sizes.  Dropping the
maximality requirement, then the disjointness restriction, and using
\(1+x\le e^x\), gives
\begin{align}\label{rt:environment-expansion}
 &\mathbb E_{\mathcal D}
 \left[\mathbf1_{\mathcal E}
  \exp\left(\kappa C\sum_F|F|
     \mathbf1_{\{d_\infty(F,\Gamma),d_\infty(F,\Gamma')\le a\}}
      \right)\right]\\
 &\hspace{2em}\le
 \exp\left(\sum_G(pe^{\kappa C})^{|G|}
     \mathbf1_{\{d_\infty(G,\Gamma),d_\infty(G,\Gamma')\le a\}}
      \right),\notag
\end{align}
where \(G\) ranges over nonempty range-\(a\) connected candidate
sets.  This argument needs only the joint bound \eqref{rt:bad-law}.

A candidate \(G\) of size \(m\) has a root within fixed distance
of some \(u_i\), with at most \(CA_a^m\) choices for the rooted
set.  If it also approaches \(\Gamma'\), a spanning tree in \(G\)
gives an opposing nominal vertex \(u'_j\) within distance \(Cm\)
of \(u_i\).  Equation~\eqref{rt:alignment} implies \(d_i\le C_{\rm rt}m\)
for a fixed integer \(C_{\rm rt}\).  Thus, increasing fixed constants,
\begin{equation}\label{rt:potential}
 \mathbb E_{\mathcal D}
  [\mathbf1_{\mathcal E}e^{\kappa J}\mid u,u']
 \le \exp\left(\sum_{i=0}^nU(d_i)\right),\qquad
 U(d)=C\kappa\mathbf1_{\{d\le L\}}
           +\sum_{m\ge1}\theta^m\mathbf1_{\{d\le C_{\rm rt}m\}},
 \quad \theta=Ap e^{C\kappa}.
\end{equation}
Here conditioning means that the nominal walks are fixed before the
environment is sampled; their independence makes its law unchanged.

Next condition on the entire opposing walk \(u'\).  In
\eqref{rt:nominal}, the random part of \(u_i\) is \(X_{s_n(i)}\).
The test \(d_i\le l\) implies that this variable belongs to a square
of radius \(l\) with deterministic transverse center
\[
 (u'_{t(i)}-z)^\perp
       -\left(\left\lfloor i\Delta_2/n\right\rfloor,
               \left\lfloor i\Delta_3/n\right\rfloor\right)
       \pmod K.
\]
The first-coordinate condition may be dropped for an upper bound.
Each folded index \(s\) occurs at most twice.  Let \(A_s(l)\) be
the union of these at most two target squares.  Then
\[
 \sum_{i=0}^n U(d_i)
 \le \sum_{s=0}^{\lfloor n/2\rfloor}V_s,\qquad
 V_s=2C\kappa\mathbf1_{\{X_s\in A_s(L)\}}
          +2\sum_{m\ge1}\theta^m
                    \mathbf1_{\{X_s\in A_s(C_{\rm rt}m)\}}.
\]
The deterministic bound
\(0\le V_s\le V_{\max}=2C\kappa+2\theta/(1-\theta)\)
implies \(e^{V_s}-1\le e^{V_{\max}}V_s\).  Choose first
\(\kappa>0\), then \(p_0>0\), sufficiently small that
\[
 2C_1e^{V_{\max}}\left[
 C\kappa(1+L)^2+\sum_{m\ge1}\theta^m(1+C_{\rm rt}m)^2\right]\le\frac12.
\]
Lemma~\ref{rt:ordered} now bounds the expectation of
\(\prod_s e^{V_s}\) by two, uniformly in the fixed opposing walk
and the two endpoints.  Averaging that walk proves
\eqref{rt:exponential}.  This conditioning argument is why equality
of the terminal endpoints is unnecessary.

\subsection{A tail bound after the silent portion}

The preceding exponential estimate controls all encounters, including
the unavoidable meeting at \(z\).  To obtain a small probability after
the silent portions, we keep track of the root index in the component
count instead of summing it by the length of the route.

We prove the following more flexible bound.  Suppose contacts are
counted only when both final-route vertices have distance at least
\(H\) from \(z\), where \(2\le H\le cR\), and denote their count
by \(J_H\).  With the site average of Proposition~\ref{rt:routes},
\begin{equation}\label{rt:general-tail}
 \mathbb E_{\rm sites}\mathbb E_{\mathcal D,\rm routes}
  [\mathbf1_{\mathcal E}\mathbf1_{\{J_H>0\}}]
 \le CC_\omega\left(H^{1-2\sigma}+\frac{H^3}{R^2}+e^{-cH}\right).
\end{equation}
Increasing constants covers bounded \(H\).  The expectation on the
left is taken successively over deterministic site weights, the datum,
and the fresh route samples; the notation below combines independent
nominal randomness and environment averages when convenient.

Increase the fixed constant \(C_{\rm rt}\) if necessary, and for \(m\ge0\) define
\begin{equation}\label{rt:root-count}
 N_m(H)=\sum_{i=0}^n
  \mathbf1_{\{d_\infty(u_i,z)\ge H-C_{\rm rt}(m+1)\}}
  \mathbf1_{\{d_\infty^\perp(u_i,u'_{t(i)})\le C_{\rm rt}(m+1)\}}.
\end{equation}
Here \(d_\infty^\perp\) uses only the last two periodic coordinates.
A contact between two nominal refinement vertices contributes to
\(N_0(H)\), after adjusting \(C_{\rm rt}\).

For a contact involving a shell, the component \(F\) used in
\eqref{rt:detour-charge} has size \(m\ge1\).  It has a root within
fixed distance of a first-route nominal vertex \(u_i\), and the
contact vertex on that route is within \(C(m+1)\) of \(u_i\).
To see the latter assertion, if that contact vertex is nominal choose
its own step; otherwise its shell-generating component was met by the
original first refinement, and a path in \(F\) joins that visit to a
bad neighbor of the contact vertex.  In either case the same \(F\)
also approaches the opposing refinement.  The diameter bound and
\eqref{rt:alignment} give both tests in \eqref{rt:root-count}.
There are at most \(A^m\) candidate rooted sets for each \(i\).
With nominal walks fixed, \eqref{rt:bad-law} therefore gives the
union bound
\begin{equation}\label{rt:tail-expansion}
 \mathbb E_{\rm sites}\mathbb E
       [\mathbf1_{\mathcal E}\mathbf1_{\{J_H>0\}}]
 \le \mathbb E_{\rm sites}\mathbb E N_0(H)
       +\sum_{m\ge1}(Ap)^m
                    \mathbb E_{\rm sites}\mathbb E N_m(H).
\end{equation}
The sum over route indices remains inside \(N_m\); in particular
this estimate has no preliminary factor depending on \(K\).

First drop the distance-from-\(z\) test in \eqref{rt:root-count}.
Condition on the opposing nominal walk and use
\eqref{rt:target-atoms} from time zero.  With the trivial bound at
folded index zero, each summand has probability at most
\(C(m+1)^2(1+s_n(i))^{-2\sigma}\).  Every folded index occurs at
most twice, so
\begin{equation}\label{rt:root-global}
 \mathbb E N_m(H)\le C(m+1)^2
\end{equation}
uniformly in both endpoints.  For a small fixed \(\delta_1>0\),
the terms of \eqref{rt:tail-expansion} with \(m>\delta_1H\) are
therefore bounded by \(Ce^{-cH}\).

Consider \(m\le\delta_1H\) and put
\(h_0=\lfloor\delta_2H\rfloor\), for a second small fixed constant
\(\delta_2>0\).  The indices with \(s_n(i)\ge h_0\) contribute at
most
\begin{equation}\label{rt:bulk-tail}
 C(m+1)^2\sum_{s\ge h_0}(1+s)^{-2\sigma}
 \le C(m+1)^2H^{1-2\sigma}.
\end{equation}
Those with \(i<h_0\) contribute zero.  Indeed bounded speed gives
\(d_\infty(u_i,z)\le3i\), whereas the distance test requires
\(d_\infty(u_i,z)\ge H-C_{\rm rt}(m+1)\).  Choose \(\delta_1,\delta_2\)
so that \(C_{\rm rt}\delta_1+3\delta_2<1/2\), and increase the lower
threshold for \(H\) if necessary.

It remains to average the indices \(i=n-s\), \(0\le s<h_0\),
near the noncommon endpoint \(y\).  Freeze \(y'\) and the entire
opposing nominal walk.  For a fixed first coordinate of \(y\), both
\(n\) and \(k_0=\min(n,n')\) are fixed.  Since clamping is
one-Lipschitz, \(|t(n-s)-k_0|\le s\).  Also
\(d_\infty(u_{n-s},y)\le3s\).  The transverse approach test in
\eqref{rt:root-count} therefore implies
\begin{equation}\label{rt:endpoint-test}
 d_\infty^\perp(y,u'_{k_0})\le6s+C_{\rm rt}(m+1).
\end{equation}
Its center is independent of the transverse coordinates of \(y\).
For each first-coordinate choice there are at most
\(C(s+m+1)^2\) possible transverse endpoint sites.  There are
\(O(R)\) first-coordinate choices in \(Q\).

To apply the joint site-weight bound without assuming a conditional
probability bound, use \(\omega(y,y')\le C_\omega R^{-6}\)
directly, and sum first in \(y\), then in the at most \(CR^3\)
choices of \(y'\).  Inequality~\eqref{rt:endpoint-test} is a
necessary condition for every first-route randomization, so its law
need not be compared across different endpoint choices.  It follows
that the site average of the remaining part of \(N_m(H)\) is at most
\begin{equation}\label{rt:endpoint-tail}
 \frac{CC_\omega}{R^2}
       \sum_{0\le s<h_0}(s+m+1)^2
 \le CC_\omega\frac{H^3}{R^2},\qquad m\le\delta_1H.
\end{equation}
Combining \eqref{rt:root-global}, \eqref{rt:bulk-tail}, and
\eqref{rt:endpoint-tail} in \eqref{rt:tail-expansion}, and using
\(\sum_m(Ap)^m(m+1)^2<\infty\), proves
\eqref{rt:general-tail}.

Every vertex used in \(J_*\) has distance greater than \(h-q\)
from \(z\).  Apply \eqref{rt:general-tail} with \(H=h-q\), adjusting
the fixed lower threshold on \(R\).  Since \(h\asymp R^{1/2}\),
\[
 H^{1-2\sigma}+H^3/R^2+e^{-cH}
 \le C\bigl(R^{-(\sigma-1/2)}+R^{-1/2}\bigr)
 \le CR^{-c_g}.
\]
This proves \eqref{rt:rare}.

Finally, viewed from the terminal endpoint, a nominal walk has the form
\[
 u_{n-i}=y+\left(-i,
       \left\lfloor\frac{-i\Delta_2}{n}\right\rfloor
                         +S^{(2)}_{s_n(i)},
       \left\lfloor\frac{-i\Delta_3}{n}\right\rfloor
                         +S^{(3)}_{s_n(i)}\right).
\]
The integer identity
\(\lfloor(n-i)\Delta_a/n\rfloor-\Delta_a
 =\lfloor-i\Delta_a/n\rfloor\) verifies this formula.  The same
two independent processes, bounded speed, clamped alignment, and
component charges apply.  This proves the common-terminal-end version
and completes Proposition~\ref{rt:routes}.

\begin{remark}\label{rt:averaging}
Both \eqref{rt:exponential} and \eqref{rt:rare} average over the
environment.  They do not assert the corresponding estimates for each
fixed datum.  They remain valid with any further retention event
inserted on their left sides.  In applications where endpoint weights
depend on the datum, a deterministic pointwise bound by
\(CR^{-6}\) permits replacing those weights before the environment
average; the counting proof of \eqref{rt:endpoint-tail} then applies
unchanged.
\end{remark}
